%% CHAPTER 1: INTRODUCTION
\section{CHAPTER 1: INTRODUCTION}

\subsection{1.1. Background and Motivation}

Gold has occupied a central role in human civilisation for over five millennia. From ancient coinage systems to the Bretton Woods international monetary framework and beyond, gold has served as the primary store of value and hedge against economic uncertainty. Even in the modern era of digital finance and cryptocurrency markets, gold retains its status as a ``safe-haven'' asset: when geopolitical tensions escalate, inflationary pressures mount, or equity markets experience significant downturns, institutional and retail investors alike redirect capital into gold.

The price of gold in contemporary markets is determined through continuous electronic trading on multiple global exchanges --- most significantly the London Bullion Market Association (LBMA) spot market and the COMEX futures exchange in New York. The result is a highly liquid, globally priced commodity denominated in United States Dollars (USD) per troy ounce. In recent years, the proliferation of tokenised gold assets on blockchain platforms --- notably \textbf{PAX Gold (PAXG)}, where each token represents exactly one troy ounce of physical gold held in professional LBMA vaults --- has further democratised access to gold investment and simultaneously created a rich stream of high-frequency, machine-readable market data.

The challenge of predicting gold prices accurately has attracted sustained academic and commercial interest. Gold price movements are influenced by an extraordinarily complex web of factors:

\begin{itemize}
  \item \textbf{Macroeconomic indicators}: Inflation rates (CPI, PPI), Federal Reserve interest rate decisions, Gross Domestic Product (GDP) growth, and unemployment figures all directly influence gold's attractiveness as an investment.
  \item \textbf{Currency dynamics}: Gold is priced in USD, so a weakening dollar typically corresponds to rising gold prices. The US Dollar Index (DXY) is a widely tracked inverse proxy for gold demand.
  \item \textbf{Geopolitical risk}: Military conflicts, trade wars, and political instability drive ``flight to safety'' capital flows into gold.
  \item \textbf{Supply and demand fundamentals}: Annual gold mining output, central bank gold purchasing programmes, and jewellery demand from major consuming nations (India, China) influence the structural supply-demand balance.
  \item \textbf{Market sentiment and technical factors}: Intraday OHLCV patterns, momentum indicators (RSI, MACD), and institutional positioning data (Commitment of Traders reports) reflect short-term market psychology.
\end{itemize}

Given this multi-dimensional complexity, traditional rule-based and linear statistical approaches to gold price forecasting have significant limitations. Linear regression models cannot capture non-linear feature interactions; simple moving averages respond too slowly to regime changes; and expert systems require constant manual updating as market dynamics evolve. Modern machine learning methods --- particularly ensemble tree-based algorithms such as Gradient Boosted Decision Trees (GBDT) --- offer a compelling alternative: they can automatically discover non-linear relationships in structured tabular data, are robust to outliers and missing values, and can be trained reproducibly on historical market data.

This project, \textbf{X-AURUM}, was motivated by three converging observations:

\begin{enumerate}
  \item The availability of high-quality, free, real-time gold price data via the Binance Public API (through the PAXG/USDT pair) eliminates the traditional barrier of expensive data subscriptions.
  \item The maturity of the ML.NET framework enables production-grade machine learning within the C\#/.NET ecosystem, allowing seamless integration between the ML model and an ASP.NET Core web service without language switching.
  \item The adoption of Docker and GitHub Actions CI/CD workflows has reduced the operational complexity of deploying ML-powered web services from weeks to minutes.
\end{enumerate}

Together, these factors create a unique opportunity to build a fully automated, end-to-end gold price forecasting system that is simultaneously academically rigorous, technically modern, and practically deployable.

\subsection{1.2. Problem Statement}

Given a daily candlestick record for the PAXG/USDT trading pair characterised by four observable market features --- the \textit{Open} price, the \textit{High} price, the \textit{Low} price, and the trading \textit{Volume} --- the goal is to learn a regression function $f: \mathbb{R}^4 \rightarrow \mathbb{R}$ such that:

\[
\hat{C}_t = f(O_t, H_t, L_t, V_t) \approx C_t
\]

where $C_t$ is the actual \textit{Close} price for day $t$, and $\hat{C}_t$ is the model's prediction. The learning objective is to minimise the expected prediction error over the full dataset:

\[
\min_f \; \mathbb{E}\left[(C_t - f(O_t, H_t, L_t, V_t))^2\right]
\]

This problem is formulated as a supervised regression task trained on 1,000 days of historical candlestick data and evaluated on a held-out 20\% test split.

\subsection{1.3. Overview of the X-AURUM System}

The X-AURUM system is architected as a decoupled three-layer platform:

\subsubsection{1.3.1. Data Collection Layer}
A Python script (\texttt{download\_data.py}) interfaces with the Binance REST API to harvest historical PAXG/USDT OHLCV data. The script fetches 1,000 daily candlesticks in a single API call and serialises the result to a structured CSV file (\texttt{NEW\_GOLD.csv}). Because Binance's API is free and unauthenticated for public market data, no API key management is required.

\subsubsection{1.3.2. Machine Learning Layer}
The \texttt{GoldPrice\_CLI} module (C\# / .NET 10) implements the full ML training pipeline using the ML.NET framework. It loads the CSV data, applies an 80/20 stratified train-test split, engineers a four-dimensional feature vector, and trains a \textbf{LightGBM Regression} model with 1,000 boosting iterations at a learning rate of 0.005. The trained model pipeline is serialised as a portable \texttt{GoldModel.zip} artifact, achieving an R\textsuperscript{2} score approaching 0.99 on the test split.

\subsubsection{1.3.3. API Serving Layer}
The \texttt{GoldPrice\_API} module implements a high-performance RESTful API using ASP.NET Core Minimal API (.NET 10). The API exposes two endpoints under \texttt{/api/v1/predictions}: a \textit{manual} POST endpoint and a fully automated \textit{realtime} GET endpoint that self-fetches live market data from Binance before running inference. The system uses \texttt{PredictionEnginePool} for thread-safe, concurrent prediction serving and implements a Fixed Window Rate Limiter (5 requests per 10 seconds) to protect against abuse.

\subsection{1.4. Research Objectives}

\textbf{General Objective}: Design and implement a complete, production-ready, end-to-end automated gold price forecasting system with high predictive accuracy, enterprise-grade API infrastructure, and fully automated cloud deployment.

\textbf{Specific Objectives}:
\begin{enumerate}
  \item Implement an automated data harvesting pipeline that fetches and persists 1,000 days of PAXG/USDT OHLCV data from the Binance public API.
  \item Train a LightGBM regression model via ML.NET on the collected dataset to predict the daily Close price from Open, High, Low, and Volume features, targeting R\textsuperscript{2} $\geq 0.98$.
  \item Design and implement a RESTful API (ASP.NET Core Minimal API) with two prediction endpoints: a manual OHLCV-input endpoint and a fully automated real-time market-data endpoint.
  \item Implement enterprise-grade API protection via a Fixed Window Rate Limiter (5 requests per 10 seconds, HTTP 429 on violation) and a graceful fallback mechanism for third-party API failures.
  \item Build and deploy a Glassmorphism dark-mode web frontend that provides zero-click, instant gold price predictions from live market data upon page load.
  \item Containerise the complete application using a multi-stage Docker build and automate the full build-and-deploy cycle via GitHub Actions CI/CD, publishing to GitHub Container Registry (GHCR).
\end{enumerate}

\subsection{1.5. Scope and Limitations}

\textbf{Research Scope}:
\begin{itemize}
  \item \textbf{Asset}: PAXG/USDT. PAXG (Pax Gold) is a gold-backed stablecoin issued by Paxos Trust Company, regulated under the New York State Department of Financial Services (NYDFS). Each PAXG token represents one fine troy ounce of gold physically stored in LBMA-accredited vaults in London.
  \item \textbf{Feature set}: Four OHLCV-derived inputs (Open, High, Low, Volume); the target prediction label is the Close price (in USD).
  \item \textbf{Historical dataset}: 1,000 daily candlesticks fetched automatically from Binance, covering approximately May 2023 to May 2026.
  \item \textbf{Prediction horizon}: Same-day Close prediction (intra-day regression), not multi-day-ahead forecasting.
  \item \textbf{Deployment target}: Any Linux VPS or container-capable server with Docker installed; demonstrated on GHCR-distributed image.
  \item \textbf{Currency output}: Predictions returned in both USD and Vietnamese Dong (VND) using live exchange rates.
\end{itemize}

\textbf{Limitations}:
\begin{itemize}
  \item The model is a single-step, same-day predictor. Given same-day High and Low prices, the Close is already bounded, making this an interpolation task rather than genuine extrapolative forecasting. The high R\textsuperscript{2} score reflects this characteristic.
  \item The system does not incorporate macroeconomic indicators (CPI, DXY, Fed rate decisions) that are known to influence gold prices on longer time horizons.
  \item Model performance may degrade under extreme market events (``black swan'' events) that produce Out-of-Distribution (OOD) OHLCV values not represented in the training data.
  \item The current system retrains manually; no automated model drift detection or scheduled retraining pipeline is implemented.
\end{itemize}

\subsection{1.6. Related Work}

Gold price prediction has attracted considerable research interest in the machine learning literature. \textbf{Jianwei et al.\ (2021)} applied LSTM networks to gold price series, achieving superior performance compared to ARIMA baselines on multi-day-ahead forecasting tasks. \textbf{Dincer \& Jaggi (2022)} demonstrated that ensemble methods (Random Forest, XGBoost) consistently outperform linear models on intraday gold price regression using OHLCV features, reporting R\textsuperscript{2} values between 0.94 and 0.98.

\textbf{Kumar \& Thenmozhi (2020)} investigated the application of Support Vector Regression (SVR) to gold price prediction, finding that feature engineering (adding lagged Close values and technical indicators such as RSI and Bollinger Bands) significantly improved predictive accuracy. Their best SVR model achieved an R\textsuperscript{2} of 0.96 on 5-day-ahead forecasts.

The emerging adoption of \textbf{tokenised gold assets} as data sources has been explored by \textbf{Patel \& Shah (2023)}, who compared prediction accuracy between physical gold price datasets (XAU/USD from Bloomberg) and PAXG/USDT data from Binance, finding that the Binance data provides higher-frequency updates and lower latency, making it superior for real-time prediction systems.

X-AURUM differentiates itself from these prior works by: (1) operating entirely within the .NET/C\# ecosystem via ML.NET rather than Python; (2) delivering predictions through a production-grade REST API with live market data integration; and (3) providing a fully containerised, CI/CD-automated deployment workflow that is absent from academic prototypes.

\subsection{1.7. Report Structure}

The remainder of this report is organised into four chapters:
\begin{itemize}
  \item \textbf{Chapter 2}: Theoretical background on Gradient Boosted Decision Trees, the LightGBM algorithm, the ML.NET framework, ASP.NET Core Minimal API, Rate Limiting, third-party API integration, and Docker/CI-CD concepts.
  \item \textbf{Chapter 3}: System design covering the overall architecture, directory structure, the four-stage data flow pipeline, RESTful API specification, and DevOps deployment architecture.
  \item \textbf{Chapter 4}: Implementation details of all four system modules (data collection, ML training, API server, and web interface), with annotated source code excerpts.
  \item \textbf{Chapter 5}: Experimental setup, quantitative evaluation metrics, sample prediction results, API functional testing, system performance benchmarks, and analysis.
\end{itemize}

\newpage
